\documentclass[11pt,a4paper]{article}
\usepackage[utf8]{inputenc}
\usepackage[T1]{fontenc}
\usepackage[french]{babel}
\usepackage{amsmath}
\usepackage{amssymb}
\usepackage{amsthm}
\usepackage{mathtools}
\usepackage{geometry}
\usepackage{enumitem}
\usepackage{xcolor}
\usepackage{titlesec}
\usepackage{fancyhdr}
\usepackage{booktabs}
\usepackage{array}
\usepackage{longtable}
\usepackage{tcolorbox}

% Configuration de la page
\geometry{margin=2.5cm}
\pagestyle{fancy}
\fancyhf{}
\fancyhead[L]{\leftmark}
\fancyhead[R]{Exercice Khi-deux}
\fancyfoot[C]{\thepage}

% Configuration des titres
\titleformat{\section}
{\Large\bfseries\color{blue!70!black}}
{}
{0em}
{}[\titlerule]

\titleformat{\subsection}
{\large\bfseries\color{blue!50!black}}
{}
{0em}
{}

\titleformat{\subsubsection}
{\normalsize\bfseries\color{blue!30!black}}
{}
{0em}
{}

% Boîtes colorées
\newtcolorbox{astucebox}{
    colback=green!5!white,
    colframe=green!75!black,
    title=\textbf{Astuce},
    fonttitle=\bfseries
}

\newtcolorbox{rappelbox}{
    colback=blue!5!white,
    colframe=blue!75!black,
    title=\textbf{Rappel},
    fonttitle=\bfseries
}

\newtcolorbox{attentionbox}{
    colback=orange!5!white,
    colframe=orange!75!black,
    title=\textbf{Attention},
    fonttitle=\bfseries
}

% Commandes personnalisées
\newcommand{\R}{\mathbb{R}}
\newcommand{\Var}{\text{Var}}
\newcommand{\E}{\mathbb{E}}
\newcommand{\chideux}{\chi^2}

% Métadonnées
\title{Exercice Complet - Test du Khi-deux\\
\large Test d'Indépendance}
\author{AmineGR03}
\date{\today}

\begin{document}

\maketitle

\tableofcontents
\newpage

\section{Introduction}

Ce document présente un exercice complet sur le test du Khi-deux d'indépendance, avec une explication détaillée de chaque étape.

\newpage

\section{Exercice : Test d'Indépendance du Khi-deux}

\subsection{Énoncé}

On souhaite tester s'il existe une relation entre le niveau de satisfaction des clients et le type de produit acheté. Un échantillon de 200 clients a été interrogé. Les résultats sont présentés dans le tableau suivant :

\begin{table}[h]
\centering
\begin{tabular}{|l|c|c|c|}
\hline
 & \textbf{Produit A} & \textbf{Produit B} & \textbf{Total} \\
\hline
\textbf{Satisfait} & 45 & 55 & 100 \\
\hline
\textbf{Non satisfait} & 35 & 65 & 100 \\
\hline
\textbf{Total} & 80 & 120 & 200 \\
\hline
\end{tabular}
\caption{Tableau de contingence : Satisfaction vs Type de produit}
\end{table}

\begin{enumerate}
    \item Formuler les hypothèses du test.
    \item Calculer les effectifs attendus sous l'hypothèse d'indépendance.
    \item Calculer la statistique du test du Khi-deux.
    \item Déterminer le nombre de degrés de liberté.
    \item Déterminer la valeur critique au seuil $\alpha = 0.05$.
    \item Conclure sur l'indépendance entre la satisfaction et le type de produit.
    \item Calculer la p-value et interpréter.
\end{enumerate}

\subsection{Rappels Théoriques}

\begin{rappelbox}
\textbf{Test du Khi-deux d'indépendance :}

On teste l'hypothèse que deux variables qualitatives sont indépendantes.

\textbf{Hypothèses :}
\begin{align}
H_0 &: \text{Les deux variables sont indépendantes} \\
H_1 &: \text{Les deux variables sont dépendantes}
\end{align}

\textbf{Statistique de test :}
$$\chideux_0 = \sum_{i=1}^{r} \sum_{j=1}^{c} \frac{(O_{ij} - E_{ij})^2}{E_{ij}}$$

où :
\begin{itemize}
    \item $O_{ij}$ : effectif observé dans la case $(i,j)$
    \item $E_{ij}$ : effectif attendu sous $H_0$ : $E_{ij} = \frac{n_i \times n_j}{n}$
    \item $r$ : nombre de lignes
    \item $c$ : nombre de colonnes
    \item $n$ : effectif total
\end{itemize}

\textbf{Degrés de liberté :}
$$\nu = (r - 1) \times (c - 1)$$

\textbf{Règle de décision :}
\begin{itemize}
    \item Si $\chideux_0 > \chideux_{\alpha, \nu}$ : on rejette $H_0$
    \item Sinon : on ne rejette pas $H_0$
\end{itemize}
\end{rappelbox}

\subsection{Solution Détaillée}

\begin{enumerate}
    \item \textbf{Formulation des hypothèses :}
    
    \begin{align}
    H_0 &: \text{La satisfaction et le type de produit sont indépendants} \\
    H_1 &: \text{La satisfaction et le type de produit sont dépendants}
    \end{align}
    
    \item \textbf{Calcul des effectifs attendus :}
    
    Sous l'hypothèse d'indépendance $H_0$, les effectifs attendus sont calculés par :
    $$E_{ij} = \frac{n_i \times n_j}{n}$$
    
    où $n_i$ est le total de la ligne $i$, $n_j$ est le total de la colonne $j$, et $n$ est l'effectif total.
    
    \begin{itemize}
        \item $E_{11}$ (Satisfait, Produit A) : $E_{11} = \frac{100 \times 80}{200} = \frac{8000}{200} = 40$
        \item $E_{12}$ (Satisfait, Produit B) : $E_{12} = \frac{100 \times 120}{200} = \frac{12000}{200} = 60$
        \item $E_{21}$ (Non satisfait, Produit A) : $E_{21} = \frac{100 \times 80}{200} = \frac{8000}{200} = 40$
        \item $E_{22}$ (Non satisfait, Produit B) : $E_{22} = \frac{100 \times 120}{200} = \frac{12000}{200} = 60$
    \end{itemize}
    
    \begin{table}[h]
    \centering
    \begin{tabular}{|l|c|c|}
    \hline
     & \textbf{Produit A} & \textbf{Produit B} \\
    \hline
    \textbf{Satisfait} & 40 & 60 \\
    \hline
    \textbf{Non satisfait} & 40 & 60 \\
    \hline
    \end{tabular}
    \caption{Effectifs attendus sous $H_0$}
    \end{table}
    
    \item \textbf{Calcul de la statistique du test :}
    
    $$\chideux_0 = \sum_{i=1}^{2} \sum_{j=1}^{2} \frac{(O_{ij} - E_{ij})^2}{E_{ij}}$$
    
    Calculons chaque terme :
    \begin{align}
    \frac{(O_{11} - E_{11})^2}{E_{11}} &= \frac{(45 - 40)^2}{40} = \frac{25}{40} = 0.625 \\
    \frac{(O_{12} - E_{12})^2}{E_{12}} &= \frac{(55 - 60)^2}{60} = \frac{25}{60} = 0.417 \\
    \frac{(O_{21} - E_{21})^2}{E_{21}} &= \frac{(35 - 40)^2}{40} = \frac{25}{40} = 0.625 \\
    \frac{(O_{22} - E_{22})^2}{E_{22}} &= \frac{(65 - 60)^2}{60} = \frac{25}{60} = 0.417
    \end{align}
    
    Donc :
    $$\chideux_0 = 0.625 + 0.417 + 0.625 + 0.417 = 2.084$$
    
    \item \textbf{Nombre de degrés de liberté :}
    
    $$\nu = (r - 1) \times (c - 1) = (2 - 1) \times (2 - 1) = 1 \times 1 = 1$$
    
    \item \textbf{Valeur critique :}
    
    Au seuil $\alpha = 0.05$ avec $\nu = 1$ degré de liberté, on consulte la table du Khi-deux :
    $$\chideux_{0.05, 1} = 3.841$$
    
    \item \textbf{Conclusion :}
    
    Puisque $\chideux_0 = 2.084 < \chideux_{0.05, 1} = 3.841$, on \textbf{ne rejette pas $H_0$} au seuil de 5\%.
    
    On conclut qu'il n'y a \textbf{pas de preuve statistique} d'une dépendance entre la satisfaction des clients et le type de produit acheté.
    
    \item \textbf{Calcul de la p-value :}
    
    La p-value est la probabilité d'observer une valeur au moins aussi extrême que celle observée sous $H_0$ :
    $$p = P(\chideux_1 \geq 2.084)$$
    
    En consultant la table du Khi-deux avec $\nu = 1$ :
    \begin{itemize}
        \item $P(\chideux_1 \geq 2.706) = 0.10$
        \item $P(\chideux_1 \geq 3.841) = 0.05$
    \end{itemize}
    
    Par interpolation, on trouve : $p \approx 0.15$
    
    Puisque $p = 0.15 > \alpha = 0.05$, on confirme qu'on ne rejette pas $H_0$.
    
    \textbf{Interprétation :} La probabilité d'observer un écart aussi grand (ou plus grand) entre les effectifs observés et attendus, si les variables étaient indépendantes, est d'environ 15\%. C'est une probabilité assez élevée, donc on ne peut pas conclure à une dépendance.
\end{enumerate}

\subsection{Tableau Récapitulatif des Calculs}

\begin{table}[h]
\centering
\small
\begin{tabular}{|c|c|c|c|}
\hline
\textbf{Case} & \textbf{Observé $O_{ij}$} & \textbf{Attendu $E_{ij}$} & \textbf{Contribution $\frac{(O_{ij}-E_{ij})^2}{E_{ij}}$} \\
\hline
(Satisfait, A) & 45 & 40 & 0.625 \\
\hline
(Satisfait, B) & 55 & 60 & 0.417 \\
\hline
(Non satisfait, A) & 35 & 40 & 0.625 \\
\hline
(Non satisfait, B) & 65 & 60 & 0.417 \\
\hline
\textbf{Total} & 200 & 200 & \textbf{2.084} \\
\hline
\end{tabular}
\caption{Détail des calculs pour le test du Khi-deux}
\end{table}

\begin{attentionbox}
\textbf{Conditions d'application du test :}
\begin{itemize}
    \item Les effectifs attendus doivent être $\geq 5$ dans au moins 80\% des cases
    \item Aucun effectif attendu ne doit être $< 1$
    \item Dans notre cas : tous les effectifs attendus sont $\geq 5$, donc le test est valide
\end{itemize}
\end{attentionbox}

\begin{astucebox}
\textbf{Méthode rapide pour un tableau 2×2 :}
Pour un tableau 2×2, on peut utiliser la formule simplifiée :
$$\chideux_0 = \frac{n(ad - bc)^2}{(a+b)(c+d)(a+c)(b+d)}$$
où $a, b, c, d$ sont les effectifs observés dans les 4 cases.
\end{astucebox}

\newpage

\section{Résumé des Étapes}

\begin{enumerate}
    \item \textbf{Formuler les hypothèses} : $H_0$ (indépendance) vs $H_1$ (dépendance)
    \item \textbf{Calculer les effectifs attendus} : $E_{ij} = \frac{n_i \times n_j}{n}$
    \item \textbf{Calculer la statistique} : $\chideux_0 = \sum \frac{(O_{ij} - E_{ij})^2}{E_{ij}}$
    \item \textbf{Déterminer les degrés de liberté} : $\nu = (r-1)(c-1)$
    \item \textbf{Comparer à la valeur critique} : $\chideux_0$ vs $\chideux_{\alpha, \nu}$
    \item \textbf{Conclure} : Rejeter ou ne pas rejeter $H_0$
    \item \textbf{Calculer la p-value} pour une interprétation plus fine
\end{enumerate}

\vspace{2cm}

\begin{center}
\textbf{Fin du document}
\end{center}

\end{document}

